% Part: normal-modal-logic
% Chapter: tableaux
% Section: introduction

\documentclass[../../../include/open-logic-section]{subfiles}

\begin{document}

\olfileid{nml}{tab}{int}

\olsection{مقدمه}

!!^{tableau}ها درخت‌هایی خاص (با انشعاب رو به پایین) از !!{signed
  formula}ها هستند؛ یعنی زوج‌هایی متشکل از یک نشانهٔ ارزش صدق
($\True$ یا~$\False$) و !!a{sentence}:
\[
\sFmla{\True}{!A} \text{ یا } \sFmla{\False}{!A}.
\]
!!^a{tableau} با شماری \emph{فرض} آغاز می‌شود. هر !!{signed formula}
بعدی با اعمال یکی از قواعد استنتاج تولید می‌شود. برخی قواعد استنتاج
یک یا چند !!{signed formula} را به انتهای درخت می‌افزایند؛ برخی دیگر
دو انتهای تازه می‌افزایند و دو شاخه پدید می‌آورند. قواعد به
!!{signed formula}هایی می‌انجامند که !!{formula} آن‌ها از !!{signed
formula}ای که قاعده بر آن اعمال شده است کم‌پیچیده‌تر است. وقتی شاخه‌ای
هم~$\sFmla{\True}{!A}$ و هم~$\sFmla{\False}{!A}$ را شامل شود، می‌گوییم
شاخه \emph{بسته} است. اگر همهٔ شاخه‌های !!a{tableau} بسته باشند، کل
!!{tableau} بسته است. !!{tableau} بسته، !!a{derivation}ای است که نشان
می‌دهد مجموعهٔ !!{signed formula}هایی که برای آغاز !!{tableau} به کار
رفته‌اند ارضاناپذیر است. می‌توان از این امر برای تعریف رابطهٔ
$\Proves$ استفاده کرد: $\Gamma \Proves !A$ اگر و تنها اگر مجموعهٔ
متناهی~$\Gamma_0 = \{!B_1, \dots, !B_n\} \subseteq \Gamma$ای وجود
داشته باشد چنان‌که !!{tableau} بسته‌ای برای فرض‌های زیر وجود داشته باشد:
\[
\{\sFmla{\False}{!A}, \sFmla{\True}{!B_1}, \dots, \sFmla{\True}{!B_n}\}.
\]

برای منطق‌های موجهاتی، هم باید مفهوم !!{signed formula} را گسترش دهیم
و هم قواعدی بیفزاییم که
\iftag{prvBox}{$\Box$\iftag{prvDiamond}{ و~$\Diamond$}{}}{$\Diamond$}
را پوشش دهند. افزون بر یک نشانه ($\True$ یا~$\False$)، !!{formula}ها
در !!{tableau}های موجهاتی \emph{پیشوندهایی}~$\sigma$ نیز دارند.
پیشوندها دنباله‌های ناتهی از اعداد صحیح مثبت‌اند؛ یعنی
$\sigma \in (\PosInt)^* \setminus \{\emptyseq\}$. هنگام نوشتن چنین
پیشوندهایی، $\tuple{\ }$ پیرامونشان را حذف می‌کنیم و !!{element}های
منفرد را به‌جای ویرگول با نقطه از هم جدا می‌کنیم. اگر~$\sigma$ یک
پیشوند باشد، آنگاه~$\sigma.n$ همان~$\sigma \concat \tuple{n}$ است؛
برای نمونه، اگر~$\sigma = 1.2.1$، آنگاه~$\sigma.3$ برابر~$1.2.1.3$
است. پس برای نمونه:
\[
\sFmla{\True}{\Box !A \lif !A}[1.2]
\]
یک \emph{!!{signed formula} پیشونددار} است (یا به‌اختصار یک
\emph{!!{formula} پیشونددار}).

به‌طور شهودی، پیشوند نام جهانی در مدلی است که ممکن است !!{formula}های
روی شاخه‌ای از !!a{tableau} را ارضا کند؛ و اگر~$\sigma$ نام جهانی
باشد، آنگاه~$\sigma.n$ نام جهانی دسترسی‌پذیر از (جهان موسوم به)
$\sigma$ است.

\end{document}
